\maketitle
\begin{abstract}
We answer three explicit continuation questions in the current ProveIt
polylogarithm programme. A centered form of Dougall's very-well-poised
summation has only even inverse powers in its asymptotic expansion. Its
Hurwitz-subtracted endpoint admits a closed formula at every nonpositive
integer resonance, and every spectral derivative has a finite expression
in Gamma derivatives, Hurwitz-zeta derivatives, and generalized Stieltjes
constants. Reciprocal-Gamma zeros are handled before coefficient extraction.
A quartic central-binomial specialization gives an all-resonance family
with explicit harmonic constants; a weighted primitive reduces the
hypergeometric order and transports the subtraction terms to colored
polylogarithms.

A second theorem gives the exact change of a one-sided Hadamard finite
part under any analytic orientation-preserving local coordinate change.
The correction is a finite residue, obeys a composition law, and specializes
to all argument derivatives of all generalized Stieltjes functions.
For a coordinate change tangent to the identity, the correction for
$\gamma_n^{(r)}$ vanishes for $n\ge 2r$, and this threshold is sharp.
Finally, a finite multiplication formula reduces three unequal integer
frequencies to ordinary colored Tornheim kernels, with explicit lower
Stieltjes terms accounting for the fixed endpoint coordinate. It also
produces convergent log-Gamma correlations and negative-order reduction
identities. Every new assertion is supplied with a proof and accompanied
by reproducible finite algebra or independent numerical diagnostics.
\end{abstract}

\noindent\textbf{Status.} The contributions are relative to the inspected
repository and incoming reports. Classical summations and product-integral
methods are credited explicitly. Historical priority for every derived
special case is not asserted. The current Gaussian $S_6$ and $S_8$
identities remain conjectural.
\clearpage
\tableofcontents
\clearpage

\section{Research scope, source baseline, and conventions}
\label{sec:scope}

\subsection{The questions addressed}
The repository's consolidated manuscript and its incoming research reports
already contain an extensive calculus of polylogarithms, harmonic sums,
Hurwitz-zeta jets, and regularized Stieltjes products. A useful continuation
must retain that accumulated proof status. The baseline used here is the
commit
\begin{center}\small\pin.\end{center}
The two requested source locations are the manuscript under
\src{Analysis/Polylogarithms/docs/manuscript} and the archives under
\src{docs/incoming} \cite{repo}.

The present work concentrates on three questions explicitly posed in the
newest incoming material:
\begin{enumerate}
\item The Gauss--Hurwitz report asks for a fixed Dougall summation with a
complete resonant coefficient law and finite Stieltjes closure
\cite[Section 9, Question 1]{incoming-gh}.
\item The Mellin--dilation and Stieltjes--harmonic reports ask for the exact
nonlinear coordinate correction to finite parts, including compatibility
with composition \cite{incoming-md,incoming-sh}.
\item The Stieltjes--harmonic report asks for three unequal integer
frequencies and a finite colored-kernel description with the local
normalization included \cite[Further research: three unequal integer
frequencies]{incoming-sh}.
\end{enumerate}
The answer to each is constructive. The first uses even Bernoulli
coefficients and a Gamma-residue identity. The second uses a local residue
that depends on a finite coordinate jet. The third uses a finite Hurwitz
multiplication formula before taking any finite part. In the third case,
integer weights do not require a new transcendental kernel: finitely many
ordinary colored Tornheim kernels suffice.

\subsection{Reading and audit boundaries}
The manuscript directory, its editorial ledger, relevant depth,
integration, differentiation, and discovery sections, and all five incoming
archives were retrieved. The complete text of the four reports most
directly connected to the new identities was examined, and the fifth
report's scope and relevant source-status discussion were checked
\cite{incoming-nhj,incoming-continuation}.
The source manifest records paths and hashes; the five archive bytes were
checked against their Git blob identifiers. This is a targeted proof
audit, not a certification of every theorem in the 417-page source book.

The inherited results used below are rederived when the proof is short:
the six-sign-region Fourier formula, the elementary Stieltjes
multiplication law, and the local residue mechanism. The longer
independent proof of Dougall's summation is classical and cited
\cite[Eq.~16.4.9]{dlmf-hyper}. The existing affine correction is recovered
as a specialization, which checks the normalization without presenting
that earlier theorem as a fresh discovery.

\subsection{Classical antecedents}
Gamma-ratio asymptotics and generalized hypergeometric endpoint summation
are classical. B\"uhring's study of partial sums includes the zero-excess
very-well-poised constant and a quartic-binomial specialization
\cite[Section 4]{buhring}. In particular, the first quartic-binomial row
below is a recovered classical identity. The extension proved here is
the centered coefficient structure, the uniform nonpositive-integer
resonance law, and its complete spectral coefficient calculus with a
specified Hurwitz subtraction.

Hurwitz product integrals and Tornheim sums were studied systematically
by Espinosa and Moll \cite{em-hurwitz,em-tornheim}; colored Tornheim
relations have a separate literature \cite{zhao}. The Fourier and
multiplication formulas for Hurwitz zeta and the Gamma connection are
standard \cite{dlmf-hurwitz}. Our unequal-frequency result is a finite
transport theorem, with the fixed-coordinate terms written out. Its
finite expression does not imply arithmetic independence or a minimal
number of period coordinates.

Analytic regularization and Hadamard finite parts also have classical
antecedents. The local theorem below is proved directly so that its
Jacobian convention, residue sign, and composition rule can all be
checked without an unspecified regularization prescription. The
specialized Stieltjes coefficient formula and sharp high-index
invariance are the concrete additions needed by the source programme.

\subsection{Conventions}
All powers of positive real numbers use the real logarithm. The spectral
and argument derivatives are distinguished by
\begin{align}
 \zeta(1+u,a)&=\frac1u+\sum_{n\ge0}\frac{(-1)^n}{n!}\gamma_n(a)u^n,
 &\gamma_0(a)&=-\psi(a), \label{eq:laurent}\\
 \zeta^{(k)}(s,a)&=\frac{\partial^k}{\partial s^k}\zeta(s,a),
 &\gamma_n^{(r)}(a)&=\frac{\mathrm d^r}{\mathrm da^r}\gamma_n(a).
\end{align}
The same superscript on $\psi$ denotes an argument derivative. Generalized
harmonic numbers are $H_N^{(j)}=\sum_{k=1}^Nk^{-j}$, with $H_0^{(j)}=0$;
$H_N=H_N^{(1)}$. The rising factorial is $(x)_N$, and $[t^k]F$ denotes a
Taylor coefficient, not a derivative. Bell polynomials are normalized by
\[
 \exp\!\left(\sum_{j\ge1}x_j\frac{t^j}{j!}\right)
 =\sum_{k\ge0}Y_k(x_1,\ldots,x_k)\frac{t^k}{k!}.
\]
The symbol $e(x)$ means $\exp(2\pi i x)$ when introduced in the Fourier
sections. On the circle, every finite part uses the right local
coordinate at each singular point. The notation $\FP_x$ is reserved for
that prescribed constant term; it does not denote an unnormalized
analytic continuation.

The proofs justify limiting operations analytically. Exact symbolic
checks verify finite identities and their implementations. Numerical
checks compare independent representations and detect sign, scale, and
transcription errors. They are not interval certificates or substitutes
for those proofs.
